\subsection{A coefficient quotient for arbitrary support}
\label{subsec:upper-unit-quotient}

For distributive lattices of rings, transition images and the
annihilator control of transition kernels already occur in
Bandelt and Petrich \cite[Lemma~4, pp.~159--160]{BandeltPetrich}.
The following argument permits arbitrary support multiplication and
uses a bound inside each additive matrix fiber. The final lifting step
is ordinary nilpotent idempotent lifting
\cite[Lemma~10.32.7, Tag~05BU]{Stacks}.

\begin{proof}[Proof of Theorem~\ref{thm:upper-unit-quotient}]
Each test in the definition of $K_q$ is an additive homomorphism of
$x$, so $K_q$ is a subgroup. If $q\leq q'$, naturality sends its test
to the old test followed by the transition
$pqr\vee1\leq pq'r\vee1$. Thus transitions preserve $K$.
For $a\in F(s)$ and $b\in F(t)$, an arbitrary context around $axb$
is
\[
 u(axb)v=(ua)x(bv),\qquad
 ua\in F(ps),\quad bv\in F(tr).
\]
Both tests have output support $psqtr\vee1$. Hence $axb\in K_{sqt}$.
This proves the two-sided ideal property with the noncommutative
order preserved. The quotient groups inherit all products,
transitions and the unit, giving the asserted inverse semiring $T'$.

If $s\geq1$, use the unit coefficient on both sides of $x\in K_s$.
The required transition is the identity of $F(s)$, so $x=0$.
Conversely, let $H$ be a coefficient ideal vanishing at all such
supports. For $x\in H_q$, every $uxv$ belongs to $H_{pqr}$ and
transports into $H_{pqr\vee1}=0$. Therefore $H\subseteq K$.
This proves maximality. Any ideal of $F/K$ vanishing above $1$ has
an inverse image containing $K$ and vanishing there, since those
groups were unchanged. Maximality forces that inverse image to equal
$K$. There is consequently no residual ideal of the stated kind.

Fix $E\in M_n(Q)$ with $E^2=E$. Its additive fiber
\[
 A_E=\prod_{i,j}F(E_{ij})
\]
is a possibly nonunital ring: each coefficient product is transported
along $E_{ik}E_{kj}\leq E_{ij}$ and the resulting contributions are
added in the target group. The quotient ring map has kernel
$N_E=\prod_{i,j}K_{E_{ij}}$.

Put $e=E_{vv}$ and $b=E_{vw}$. Matrix idempotence gives $eb\leq b$,
and hence
\[
 (e\vee1)b=eb\vee b=b.
\]
For $z\in K_e$, the unit context makes $t_{e,e\vee1}(z)=0$.
Multiplication naturality therefore gives, for every $y\in F(b)$,
\begin{equation}\label{eq:upper-unit-cycle}
 t_{eb,b}(zy)=0.
\end{equation}
This uses the displayed matrix relation, not an assumption $e\leq1$.

Expand a product of $n+1$ matrices from $N_E$ into path
contributions. A path has vertices $i_0,\ldots,i_{n+1}$.
Among its first $n+1$ vertices two coincide, giving a nonempty cycle
with a following edge $v\to w$. The cycle coefficient belongs to $K$
at its scalar product support. Transport it to $F(E_{vv})$; this is
allowed because every positive power of $E$ is $E$, and it remains
in $K_{E_{vv}}$. Equation~\eqref{eq:upper-unit-cycle} annihilates the
cycle followed by the next edge after transport to $F(E_{vw})$.
Replacing this segment by the single $v\to w$ edge leaves a
nonempty valid path with the original endpoints. All its scalar
supports remain bounded by the final entry of $E$. Naturality
identifies the original contribution with the resulting zero
contribution. Thus $N_E^{n+1}=0$. Only the matrix index set was
required to be finite.

To lift an idempotent over $T'$, lift its entries at the same support
$E$ to $a\in A_E$. Its defect $d=a^2-a$ belongs to $N_E$.
Replace $a$ by
\[
 a'=3a^2-2a^3=a+d-2ad.
\]
The new defect is $d^2(4d-3)$, so iteration reaches an idempotent
without changing the quotient. The correction has no constant term
and remains in the possibly nonunital ring $A_E$.

For uniqueness, let lifted idempotents $P,Q$ have isomorphic quotient
images. Lift inverse rectangular maps at their exact supports, then
replace them by $X=PXQ$ and $Y=QYP$. The support equations are
unchanged. We have $XY=P+\delta$ and $YX=Q+\epsilon$, with the
defects in the respective nilpotent coefficient-fiber ideals. In
the corner rings with identities $P,Q$, both products are invertible
by finite geometric series. The identity $(YX)Y=Y(XY)$ gives
\[
 Y'=Y(XY)^{-1}=(YX)^{-1}Y,\qquad XY'=P,\quad Y'X=Q.
\]
These are the required lifted inverse maps. Coefficient zero at a
nonzero support is retained throughout; it is not the zero semiring
matrix.

Extension of scalars sends a split idempotent image to its quotient
idempotent image. The existence and uniqueness just proved therefore
give its bijection on finite-projective classes. Block sums are
preserved, yielding the stated monoid and group isomorphisms.
\end{proof}

The largest-ideal characterization makes the construction invariant
under pair isomorphisms. It does not make it functorial under an
arbitrary map to another pair: the target may contain coefficient
contexts unavailable in the source. The theorem reduces the general
coefficient diagram; it does not classify the residual projective
supports or ordinary amplitudes.
